\section{Inciso d)}

% Seccion con nombres de clase largos, que no se parten con guion.
\sloppy

\textit{Compara: Visitor vs. añadir métodos a cada elemento. ¿Cuándo conviene cada uno? Usa la
métrica de ``número de cambios por nueva operación'' y ``número de cambios por nuevo tipo de
elemento''. ¿Qué pasa si el número de operaciones crece mucho más rápido que el número de
tipos de cuenta?}

\vspace{0.4cm}

\subsection*{Las dos alternativas son la misma matriz, descompuesta de dos maneras}

El comportamiento que el sistema debe implementar es una matriz de $T \times V$ celdas: $T$
tipos de cuenta por $V$ operaciones. En esta implementación son $T = 4$ y $V = 5$, es decir
veinte comportamientos.

Las dos alternativas no se diferencian en qué pueden expresar, porque ambas llenan la misma
matriz, sino en \textbf{cómo la agrupan en unidades de código}:

\begin{itemize}[leftmargin=1.2cm]
    \item \textbf{Métodos en cada elemento:} la matriz se parte por filas. Cada clase de
          cuenta contiene sus $V$ comportamientos.
    \item \textbf{Visitor:} la matriz se parte por columnas. Cada visitante contiene los $T$
          comportamientos de una operación.
\end{itemize}

Y como una clase es la unidad de modificación, la elección decide \textbf{qué eje queda
cerrado a modificación}: el eje que se corta en unidades es el que se puede extender sin
tocar nada. Esto es el \textit{expression problem}, y no tiene solución gratuita en un
lenguaje con despacho simple: lo que se gana en un eje se paga en el otro.

\subsection*{Las dos métricas}

Para que la métrica signifique algo hay que distinguir dos cosas que no cuestan lo mismo:
\textbf{añadir} un archivo nuevo, que no puede romper lo que ya funciona, y
\textbf{modificar} un archivo existente, que sí puede. Un archivo modificado obliga a releer
el código que lo rodea, a repetir sus pruebas, a recompilar cuanto dependa de él y, en un
equipo que trabaja en ramas paralelas, es un conflicto de fusión en potencia. La métrica del
enunciado cuenta, por tanto, archivos modificados:

\begin{table}[h!]
    \centering
    \small
    \begin{tabular}{p{3.9cm}p{5.1cm}p{5.1cm}}
        \toprule
        & \textbf{Cambios por nueva operación} & \textbf{Cambios por nuevo tipo de elemento} \\
        \midrule
        \textbf{Visitor}
            & 1 archivo nuevo, \textbf{0 modificados}
            & 1 archivo nuevo y \textbf{$V+1$ modificados}: la interfaz del visitante y los $V$
              visitantes \\
        \addlinespace
        \textbf{Métodos en los elementos}
            & \textbf{$T+1$ modificados}: las $T$ clases de cuenta y su interfaz común
            & 1 archivo nuevo, \textbf{0 modificados} \\
        \bottomrule
    \end{tabular}
\end{table}

La tabla es perfectamente simétrica, y en esa simetría está todo el argumento: no hay una
opción mejor, hay una opción adecuada a cada perfil de crecimiento.

\subsection*{Medición sobre esta implementación}

Las dos cifras no se afirman, se midieron sobre el propio repositorio.

\textbf{Coste de una operación nueva con Visitor.} La quinta operación,
\texttt{ExportarCsvVisitor}, se añadió en un \textit{commit} aislado precisamente para que la
medida fuera verificable:

{\small
\begin{verbatim}
$ git show --stat --format="%h %s" da8a926
da8a926 docs(cp4): agrega una quinta operacion como experimento del inciso (d)

 .../CP4/codigo/ejercicio12/ExportarCsvVisitor.java | 40 ++++++++++++++++++++++
 1 file changed, 40 insertions(+)
\end{verbatim}
}

Un archivo nuevo, cuarenta líneas, \textbf{cero archivos modificados}. Ninguna de las cuatro
cuentas se tocó, y las cuatro operaciones anteriores siguieron compilando y pasando sus
demostraciones. El contrafactual es igual de concreto: la misma operación escrita como método
de los elementos habría modificado cinco archivos, las cuatro cuentas más
\texttt{ElementoContable}, y habría forzado la recompilación de todo lo que depende de esa
interfaz.

\textbf{Coste de un tipo de elemento nuevo con Visitor.} El inciso (c) introdujo
\texttt{GrupoCuentas}, que es exactamente eso: un tipo de elemento nuevo. El coste obligatorio
fue un archivo nuevo más \emph{una} modificación, la interfaz \texttt{VisitanteOperacion}.
Pagó menos que los $V+1$ de la tabla por una razón concreta: los métodos añadidos son
\texttt{default}, de modo que los cuatro visitantes existentes no tuvieron que cambiar. Sin
ese recurso del lenguaje habrían sido cinco archivos modificados.

\begin{table}[h!]
    \centering
    \small
    \begin{tabular}{p{5.6cm}p{2.5cm}p{2.5cm}p{3.0cm}}
        \toprule
        \textbf{Cambio realizado} & \textbf{Archivos nuevos} & \textbf{Modificados}
        & \textbf{Con métodos en los elementos} \\
        \midrule
        Quinta operación (exportar a CSV) & 1 & 0 & 5 modificados \\
        \addlinespace
        Quinto tipo de elemento (el grupo) & 1 & 1 & 0 modificados \\
        \bottomrule
    \end{tabular}
\end{table}

La comprobación de que las cinco operaciones conviven sobre la misma jerarquía intacta es la
salida de \texttt{DemoMetricas}:

{\small
\begin{verbatim}
Operacion 5 (agregada sin tocar las cuentas):
AHO-001;Ahorro;4,500.00;retiros=1
COR-001;Corriente;-305.00;sobregiro=500.00
INV-001;Inversion;19,490.00;penalizaciones=10.00
EMP-001;Empresarial;99,500.00;doble-firma=true

Las cinco operaciones sobre el mismo arbol:
  1 ejecutar   -> EMP-001
  2 revertir   -> true
  3 auditar    -> riesgo ALTO
  4 intereses  -> 1,996.10
  5 exportar   -> 4 lineas

Archivos de elemento modificados para lograrlo: 0
\end{verbatim}
}

\subsection*{Cuándo conviene cada uno}

El criterio de decisión se puede enunciar en una frase: \textbf{el eje que se cierra debe ser
el que no va a crecer}, y si alguno de los dos lo van a extender terceros, ése es el que hay
que dejar abierto.

\textbf{Conviene Visitor cuando:}

\begin{itemize}[leftmargin=1.2cm]
    \item La jerarquía de elementos es cerrada y estable, y está bajo control propio.
    \item Las operaciones son numerosas o se anuncian como crecientes, que es el caso del
          enunciado.
    \item La operación \emph{no} es responsabilidad del elemento: auditar, exportar, calcular
          intereses regulatorios o generar reportes no forman parte de lo que una cuenta
          \emph{es}.
    \item La operación necesita recorrer la estructura y agregar resultados, como la
          consolidación del inciso (c).
\end{itemize}

\textbf{Conviene poner métodos en el elemento cuando:}

\begin{itemize}[leftmargin=1.2cm]
    \item La operación es intrínseca al elemento y necesita su estado privado. En esta misma
          implementación, \texttt{saldo()}, \texttt{aplicar()} y \texttt{cerrojo()} son
          métodos de \texttt{Cuenta}, y sacarlos a un visitante sería absurdo.
    \item Las operaciones son pocas y estables: con $V$ pequeño, la indirección de Visitor
          solo resta legibilidad.
    \item La jerarquía de elementos crece, sobre todo si la extienden terceros. Una biblioteca
          cuyos usuarios definen subclases no puede cerrar ese eje.
\end{itemize}

Lo que esta implementación muestra, y conviene subrayar, es que \textbf{las dos cosas conviven
sin contradicción}: \texttt{Cuenta} tiene métodos propios para lo que es suyo, y las
operaciones de negocio viven en visitantes. La pregunta correcta ante cada operación no es qué
patrón usar, sino si esa operación define al elemento o si es una política externa que lo
consulta.

\subsection*{Qué pasa si las operaciones crecen mucho más rápido que los tipos}

Es la situación del enunciado, $V \gg T$ con $T$ fijo, y tiene cuatro consecuencias.

\textbf{1. La ventaja acumulada de Visitor es lineal en $V$.} Cada operación nueva evita
$T+1$ modificaciones, así que llegar a $V$ operaciones evita $V \cdot (T+1)$ modificaciones de
archivos existentes. Con $T = 4$ y $V = 20$, son cien modificaciones evitadas frente a cero
realizadas. En un equipo que trabaja con ramas y revisiones, ese número no es estético: son
cien ocasiones de conflicto de fusión y cien revisiones de código que ya estaba cerrado.

\textbf{2. El riesgo se concentra en el eje cerrado.} El coste de un tipo de elemento nuevo es
$V+1$, y ese número crece con el éxito del diseño: con $V = 20$, un quinto tipo de cuenta
cuesta veintiuna modificaciones. Visitor cambia un coste frecuente y moderado por uno raro y
caro, lo cual es un buen negocio \emph{mientras la premisa se cumpla}. Si la jerarquía resulta
no estar cerrada, el patrón cobra la deuda entera de golpe.

\textbf{3. Hay mitigaciones, y todas tienen precio.} Conviene conocerlas y elegir a ojos
abiertos:

\begin{itemize}[leftmargin=1.2cm]
    \item \textbf{Métodos \texttt{default} en la interfaz del visitante}, que es lo hecho en
          el inciso (c): el tipo nuevo cuesta una sola modificación. El precio es que el
          comportamiento por omisión es silencioso, y un visitante que debería tratar el tipo
          nuevo puede ignorarlo sin que el compilador diga nada. Se cambia seguridad en
          compilación por comodidad.
    \item \textbf{Un visitante base abstracto con un \texttt{visitPorDefecto}}, la variante
          conocida como \textit{Default Visitor}: el mismo compromiso, concentrado en una
          clase que los demás heredan.
    \item \textbf{Dejar la interfaz sin valores por omisión}: el compilador señala uno por uno
          los $V$ visitantes que hay que revisar. Es el camino caro y el único seguro; con $V$
          grande sigue siendo el correcto cuando el tipo nuevo cambia la semántica de las
          operaciones existentes.
\end{itemize}

\textbf{4. Aparecen costes que la métrica no mide.} Contar archivos modificados es una buena
aproximación al riesgo, pero deja fuera tres cosas que con $V$ grande se vuelven dominantes:

\begin{itemize}[leftmargin=1.2cm]
    \item \textbf{Descubribilidad.} Con cincuenta operaciones, leer \texttt{CuentaAhorro} ya
          no dice qué se puede hacer con una cuenta de ahorro: hay que buscar los cincuenta
          visitantes. La respuesta a ``qué sabe hacer este objeto'' deja de estar en el
          objeto.
    \item \textbf{Duplicación entre visitantes.} Recorridos, acumuladores y validaciones se
          repiten y hay que factorizarlos en visitantes base, o el ahorro de modificaciones se
          paga en código repetido.
    \item \textbf{Pruebas.} La matriz $T \times V$ crece con $V$: con $T = 4$ y $V = 20$ son
          ochenta combinaciones. Conviene una tabla de casos parametrizada antes que una
          prueba por método.
\end{itemize}

\subsection*{Una nota sobre la alternativa moderna}

La asimetría que obliga a elegir no es una ley del diseño orientado a objetos, sino una
consecuencia del despacho simple. Con tipos sellados y \textit{switch} exhaustivo sobre
patrones, disponibles como característica definitiva desde Java 21, una operación nueva es una
función nueva, con cero modificaciones, y un tipo nuevo hace \emph{fallar la compilación} en
todos los \textit{switch} que lo ignoran, de modo que el compilador enumera exactamente los
sitios a revisar. El compromiso sigue existiendo, pero el eje caro deja de ser silencioso.

Esta implementación usa Java 17, donde el \textit{switch} sobre patrones era aún una
característica en vista previa, así que Visitor sigue siendo la forma canónica de obtener
despacho doble con comprobación en tiempo de compilación. Vale la pena decirlo porque fija el
límite del argumento: se eligió Visitor por el perfil de crecimiento del problema y por las
herramientas del lenguaje disponible, no porque sea una verdad permanente.

\subsection*{Conclusión para el problema del banco}

Los tipos de cuenta están fijados en cuatro por el enunciado y las operaciones se declaran
explícitamente crecientes: $V \gg T$ con $T$ estable. Bajo esa condición, la comparación no
está reñida. Visitor cuesta cero modificaciones por operación nueva, frente a cinco, y la
medición sobre el repositorio confirma la cifra. El riesgo asumido queda nombrado: si
apareciera un quinto tipo de cuenta, el coste serían $V+1$ modificaciones, rebajado a una sola
por los métodos \texttt{default} del inciso (c), al precio de que el compilador deje de avisar
de los visitantes que habría que revisar.

\fussy
